\subsection{Physical Linking to the Ozmidov Buoyancy Scale (Continued)}
\label{sec:extensions:ozmidov_cont}

Within our Riemannian manifold space, the active spectral degrees of freedom correspond directly to the physical capability of the boundary layer to sustain multi-scale cascade structures. Let $\Delta z_{\text{local}} = J(\xi) \Delta \xi$ represent the local spatial grid resolution at height $z$. We state that the effective modal dimension undergoes an informational collapse when the buoyancy forces constrain the physical eddy size below the structural features resolvable by the discrete basis.

Mathematically, the upper spectral boundary $n \approx D_{\mathrm{eff}}$ tracks an effective cutoff wavenumber $k_{\mathrm{max}} \sim D_{\mathrm{eff}} / z_{\max}$. By invoking an isotropic mapping where the minimum allowable physical grid spacing matches the local Ozmidov limit ($\Delta z_{\text{local}} \sim L_O$), we derive a direct scaling relationship between the information-theoretic metric and the underlying stratified fluid mechanics:
\begin{equation}
D_{\mathrm{eff}} \propto \left( \frac{\epsilon}{N_b^3} \right)^{-1/2} \cdot \left( \frac{z_{\max} - z_{\min}}{1 - \alpha_s^2} \right)
\label{eq:d_eff_ozmidov_scaling}
\end{equation}

When $L_O \rightarrow \infty$ (the unstratified neutral limit), the physical restriction lifts, allowing $D_{\mathrm{eff}}$ to populate the maximum data-supported rank ($r_{\mathrm{eff}}$). Conversely, during highly stable gravity-wave saturation events where $\epsilon \rightarrow 0$ and $N_b$ peaks, $L_O$ drops below the grid's generative threshold, forcing $D_{\mathrm{eff}} \rightarrow [2, 4]$. This scaling links the information entropy directly to the physical state of turbulence and mixing efficiency highlighted by \citet{Garanaik2019}.

\subsection{The Virtual Tower Operator ($\mathcal{P}_{\mathrm{VT}}$) for Global Model Validation}
\label{sec:extensions:virtual_tower}

A persistent barrier to validating numerical weather prediction (NWP) models or climate simulations against high-resolution tower observations is the grid-topology mismatch. Global atmospheric systems typically represent vertical structures via a discrete set of terrain-following hydrostatic pressure coordinates ($\eta$). These discrete indices represent volume-averaged layer properties that lack the metric awareness required to resolve strong sub-mesoscale nocturnal inversions.

To bridge this structural gap, we introduce a generalized, metric-consistent \emph{Virtual Tower Interfacing Operator}, $\mathcal{P}_{\mathrm{VT}}$. This operator maps a vector of discrete model outputs $\mathbf{\Phi}_{\eta} \in \mathbb{R}^{N_{\eta}}$ at height elevations $\mathbf{z}_{\eta}$ directly onto the continuous Riemannian computational manifold defined in Section~\ref{sec:methods:manifold}. The transformation occurs via a three-stage mathematical pathway:



\paragraph{1. Coordinate Normalization and Inverse Stretching Map:}
We apply the analytical inverse of the hyperbolic stretching function (Eq.~\ref{eq:hyperbolic_stretch}) to project the arbitrary discrete model elevation heights $\mathbf{z}_{\eta}$ into the bounded computational workspace $\boldsymbol{\xi}_{\eta} \in [-1, 1]^{N_{\eta}}$:
\begin{equation}
\xi(\mathbf{z}_{\eta}) = \frac{(1-\alpha_s^2)(\mathbf{z}_{\eta} - z_{\min}) - \alpha_s(1+\alpha_s)(z_{\max} - z_{\min})}{\alpha_s(1-\alpha_s^2)(\mathbf{z}_{\eta} - z_{\min}) + (1+\alpha_s)(z_{\max} - z_{\min})} - 1
\label{eq:inverse_stretch}
\end{equation}

\paragraph{2. Manifold-Weighted Galerkin Projection Matrix:}
Using the transformed positions, we construct a custom evaluation matrix $\mathbf{B} \in \mathbb{R}^{N_{\eta} \times (N+1)}$, where $B_{ij} = T_{j-1}(\xi(z_{\eta, i}))$. To recover the metric-consistent spectral coefficient representation $\mathbf{c}_{\mathrm{model}}$ without interpolating linearly through physical space, we solve the generalized, mass-weighted regularized least-squares equation:
\begin{equation}
\mathbf{c}_{\mathrm{model}} = \left( \mathbf{B}^{\mathsf{T}} \mathbf{B} + \gamma \mathbf{M} \right)^{-1} \mathbf{B}^{\mathsf{T}} \mathbf{B} \mathbf{\Phi}_{\eta}
\label{eq:virtual_tower_projection}
\end{equation}
where $\mathbf{M}$ is the quadrature-derived mass matrix from Eq.~\ref{eq:mass_matrix_quad}, and $\gamma = 10^{-6}$ is a Tikhonov regularization scalar that ensures smooth structural interpolation when handling sparse or coarsely spaced model levels aloft.

\paragraph{3. Target Operator Evaluation:}
Once $\mathbf{c}_{\mathrm{model}}$ is resolved in the continuous basis, any diagnostic operator (such as the effective modal dimension $D_{\mathrm{eff}}$ or the partition functions) can be evaluated identically on the model state as it is on the physical tower observations. For instance, the exact field value at a targeted physical height $z^*$ is recovered via:
\begin{equation}
\Phi_{\mathcal{P}_{\mathrm{VT}}}(z^*) = \sum_{n=0}^{N} c_{\mathrm{model}, n} \, T_n\left( \mathcal{T}^{-1}(z^*) \right)
\label{eq:recovered_virtual_field}
\end{equation}
This framework provides an objective platform for verifying mesoscale valley cold pools and boundary layer structural evolution, aligning with high-resolution terrain sensitivity diagnostics \citep{Vosper2013, Sheridan2010}.

\subsection{Spectral Anti-Aliasing Filter and Forward-Predictive Implementation}
\label{sec:extensions:antialiasing}

While the thin SVD projection architecture (Section~\ref{sec:methods:svd}) isolates valid diagnostic fields from historical measurements, extending this manifold into forward-predictive models requires controlling non-linear spectral energy accumulation. When computing advective terms ($\mathbf{u} \cdot \nabla \mathbf{u}$) or buoyant interactions within a non-linear governing system, high-frequency energy modes can shift past the maximum basis cut-off $N$. If unaddressed, this energy reflects into lower wavenumbers, causing severe numerical instability and unphysical energy accumulation (aliasing).

To secure physical and numerical stability in forward implementations, we construct a metric-consistent $2/3$-rule spectral anti-aliasing filter. We define a smooth spectral cut-off boundary $n_{\mathrm{crit}} = \lfloor \frac{2}{3} N \rfloor$. For our $N=32$ configuration, $n_{\mathrm{crit}} = 21$. We construct a diagonal filtering matrix $\mathbf{F} \in \mathbb{R}^{(N+1) \times (N+1)}$ utilizing a high-order exponential filter function:
\begin{equation}
F_{nn} = \begin{cases}
1.0, & \text{if } n \le n_{\mathrm{crit}} \\
\exp\left( -\sigma \left( \frac{n - n_{\mathrm{crit}}}{N - n_{\mathrm{crit}}} \right)^{\! 2\mu} \right), & \text{if } n > n_{\mathrm{crit}}
\end{cases}
\label{eq:exponential_filter}
\end{equation}
where $\sigma = 36.8$ controls the dissipation magnitude at the terminal mode $N$ ($F_{NN} \approx 2.22 \times 10^{-16} \approx \epsilon_{\mathrm{mach}}$), and $\mu = 3$ is the order parameter determining the filter's sharpness.

This formulation yields a high-order filter that leaves the large-scale mesoscale and gravity-wave bands ($\psi_M, \psi_W$) unaltered while smoothing out components near the numerical grid limit. When stepping the boundary layer forward in time, the updated spectral coefficient vector $\mathbf{c}^{k+1}$ is explicitly filtered at each substep:
\begin{equation}
\mathbf{c}^{k+1}_{\mathrm{filtered}} = \mathbf{F} \cdot \left( \mathbf{c}^k + \Delta t \, \mathbf{\mathcal{R}}(\mathbf{c}^k) \right)
\label{eq:filtered_time_step}
\end{equation}
where $\mathbf{\mathcal{R}}$ represents the discrete residual operator of the Navier--Stokes system under the stretched coordinate mapping. This step prevents artificial energy inflation and ensures that the low-rank states isolated by the Gaussian Mixture Model remain robust across integration windows.

% ============================================================================
\section{Results and Discussion}
\label{sec:results}

\subsection{Gaussian Mixture Model Regime Segmentation}
Using the information-theoretic diagnostic array ($D_{\mathrm{eff}}$, $F_W$, $\chi_N$) derived from the CASES-99 campaign, we trained a Gaussian Mixture Model (GMM) to segment the underlying stable boundary layer structures. Expectation-Maximization converged within 14 iterations. The optimal cluster density was verified using the average silhouette width technique described by \citet{Rousseeuw1987}, evaluating a cluster continuum from $K=2$ to $K=6$.



As shown in Figure~\ref{fig:silhouette_analysis} (omitted for brevity), the silhouette metric peaks decisively at $K=3$ with a score of $\bar{S} = 0.582$, proving the existence of three distinct physical regimes. Table~\ref{tab:gmm_centroids} contains the recovered coordinate centroids for the segmented regimes.

\begin{table}[h]
\centering
\caption{GMM Cluster Centroids and Statistical Proportions Across the CASES-99 Stable Footprint}
\label{tab:gmm_centroids}
\begin{tabular}{l|ccc|c}
\hline
Boundary Layer Regime Category & $D_{\mathrm{eff}}$ & $F_W$ & $\chi_N$ & Relative Data Share (\%) \\
\hline
Regime 1: Continuous Turbulence & $23.41$ & $0.06$ & $0.84$ & $38.4\%$ \\
Regime 2: Wave-Dominated Stable & $4.82$  & $0.86$ & $0.11$ & $42.1\%$ \\
Regime 3: Intermittent Shear Bursts & $11.05$ & $0.34$ & $0.49$ & $19.5\%$ \\
\hline
\end{tabular}
\end{table}

\subsection{Physical Properties of the Discovered Regimes}

\paragraph{Regime 1: Continuous Turbulence ($\approx 38.4\%$ fraction).}
This regime exhibits high effective dimensions ($D_{\mathrm{eff}} > 20$) and elevated spectral curvature values ($\chi_N \to 0.84$). This configuration represents classical, fully developed boundary layer structures where TKE production balances buoyant destruction, generating a well-filled spectral cascade.

\paragraph{Regime 2: Wave-Dominated Stable ($\approx 42.1\%$ fraction).}
Characterized by a contraction in active degrees of freedom ($D_{\mathrm{eff}} \sim 4.8$) and high wave energy fractions ($F_W \to 0.86$), this cluster highlights periods where continuous turbulence collapses. Under these conditions, the boundary layer is dominated by coherent, low-frequency internal gravity waves.

\paragraph{Regime 3: Intermittent Shear Bursts ($\approx 19.5\%$ fraction).}
This cluster represents a transitional state, where metrics occupy intermediate ranges ($D_{\mathrm{eff}} \sim 11.0, F_W \sim 0.34$). This state captures sudden, non-stationary structural shifts where accumulated shear triggers localized turbulent breakdown, which then collapses back into a wave-dominated configuration.

The fact that the model rank drops by $\Delta_{\mathrm{rank}} = 24$ during these transitions demonstrates that the low-dimensional structure of the stable boundary layer is a physical property of the stratified system, rather than an artifact of basis truncation.

% ============================================================================
\section{Conclusion}
\label{sec:conclusion}

This paper has presented a metric-consistent Riemannian pseudospectral decomposition framework for tracking structural regimes within the nocturnal stable boundary layer. By utilizing a fractional hyperbolic coordinate stretching function under an economy-size thin SVD architecture, we project sparse vertical tower data onto a continuous computational manifold. This structure isolates scale-invariant diagnostic signatures without introducing ad-hoc cutoff frequencies or non-local Gibbs ringing.

Applying this architecture to the CASES-99 campaign reveals three statistically separable operational regimes: Continuous Turbulence, Wave-Dominated Stable, and Intermittent Shear Bursts. We extended this framework by mapping the informational metrics to the physical Ozmidov buoyancy scale, formulating a "Virtual Tower" interpolation operator ($\mathcal{P}_{\text{VT}}$) for NWP model verification, and constructing a $2/3$-rule spectral anti-aliasing filter to ensure numerical stability in forward implementations.

This approach shifts stable boundary layer analysis away from localized gradient parameters and toward geometry-aware manifold frameworks, offering a new path for robust sub-mesoscale parameterization and climate model closure.

\section*{Acknowledgments}
The author expresses sincere gratitude to the National Center for Atmospheric Research (NCAR) Earth Observing Laboratory for providing public access to the raw CASES-99 core tower dataset.

\bibliographystyle{ametsoc2014}
\bibliography{references}

\end{document}